% Options for packages loaded elsewhere
\PassOptionsToPackage{unicode}{hyperref}
\PassOptionsToPackage{hyphens}{url}
\documentclass[
]{article}
\usepackage{xcolor}
\usepackage[margin=1in]{geometry}
\usepackage{amsmath,amssymb}
\setcounter{secnumdepth}{-\maxdimen} % remove section numbering
\usepackage{iftex}
\ifPDFTeX
  \usepackage[T1]{fontenc}
  \usepackage[utf8]{inputenc}
  \usepackage{textcomp} % provide euro and other symbols
\else % if luatex or xetex
  \usepackage{unicode-math} % this also loads fontspec
  \defaultfontfeatures{Scale=MatchLowercase}
  \defaultfontfeatures[\rmfamily]{Ligatures=TeX,Scale=1}
\fi
\usepackage{lmodern}
\ifPDFTeX\else
  % xetex/luatex font selection
\fi
% Use upquote if available, for straight quotes in verbatim environments
\IfFileExists{upquote.sty}{\usepackage{upquote}}{}
\IfFileExists{microtype.sty}{% use microtype if available
  \usepackage[]{microtype}
  \UseMicrotypeSet[protrusion]{basicmath} % disable protrusion for tt fonts
}{}
\makeatletter
\@ifundefined{KOMAClassName}{% if non-KOMA class
  \IfFileExists{parskip.sty}{%
    \usepackage{parskip}
  }{% else
    \setlength{\parindent}{0pt}
    \setlength{\parskip}{6pt plus 2pt minus 1pt}}
}{% if KOMA class
  \KOMAoptions{parskip=half}}
\makeatother
\usepackage{color}
\usepackage{fancyvrb}
\newcommand{\VerbBar}{|}
\newcommand{\VERB}{\Verb[commandchars=\\\{\}]}
\DefineVerbatimEnvironment{Highlighting}{Verbatim}{commandchars=\\\{\}}
% Add ',fontsize=\small' for more characters per line
\usepackage{framed}
\definecolor{shadecolor}{RGB}{248,248,248}
\newenvironment{Shaded}{\begin{snugshade}}{\end{snugshade}}
\newcommand{\AlertTok}[1]{\textcolor[rgb]{0.94,0.16,0.16}{#1}}
\newcommand{\AnnotationTok}[1]{\textcolor[rgb]{0.56,0.35,0.01}{\textbf{\textit{#1}}}}
\newcommand{\AttributeTok}[1]{\textcolor[rgb]{0.13,0.29,0.53}{#1}}
\newcommand{\BaseNTok}[1]{\textcolor[rgb]{0.00,0.00,0.81}{#1}}
\newcommand{\BuiltInTok}[1]{#1}
\newcommand{\CharTok}[1]{\textcolor[rgb]{0.31,0.60,0.02}{#1}}
\newcommand{\CommentTok}[1]{\textcolor[rgb]{0.56,0.35,0.01}{\textit{#1}}}
\newcommand{\CommentVarTok}[1]{\textcolor[rgb]{0.56,0.35,0.01}{\textbf{\textit{#1}}}}
\newcommand{\ConstantTok}[1]{\textcolor[rgb]{0.56,0.35,0.01}{#1}}
\newcommand{\ControlFlowTok}[1]{\textcolor[rgb]{0.13,0.29,0.53}{\textbf{#1}}}
\newcommand{\DataTypeTok}[1]{\textcolor[rgb]{0.13,0.29,0.53}{#1}}
\newcommand{\DecValTok}[1]{\textcolor[rgb]{0.00,0.00,0.81}{#1}}
\newcommand{\DocumentationTok}[1]{\textcolor[rgb]{0.56,0.35,0.01}{\textbf{\textit{#1}}}}
\newcommand{\ErrorTok}[1]{\textcolor[rgb]{0.64,0.00,0.00}{\textbf{#1}}}
\newcommand{\ExtensionTok}[1]{#1}
\newcommand{\FloatTok}[1]{\textcolor[rgb]{0.00,0.00,0.81}{#1}}
\newcommand{\FunctionTok}[1]{\textcolor[rgb]{0.13,0.29,0.53}{\textbf{#1}}}
\newcommand{\ImportTok}[1]{#1}
\newcommand{\InformationTok}[1]{\textcolor[rgb]{0.56,0.35,0.01}{\textbf{\textit{#1}}}}
\newcommand{\KeywordTok}[1]{\textcolor[rgb]{0.13,0.29,0.53}{\textbf{#1}}}
\newcommand{\NormalTok}[1]{#1}
\newcommand{\OperatorTok}[1]{\textcolor[rgb]{0.81,0.36,0.00}{\textbf{#1}}}
\newcommand{\OtherTok}[1]{\textcolor[rgb]{0.56,0.35,0.01}{#1}}
\newcommand{\PreprocessorTok}[1]{\textcolor[rgb]{0.56,0.35,0.01}{\textit{#1}}}
\newcommand{\RegionMarkerTok}[1]{#1}
\newcommand{\SpecialCharTok}[1]{\textcolor[rgb]{0.81,0.36,0.00}{\textbf{#1}}}
\newcommand{\SpecialStringTok}[1]{\textcolor[rgb]{0.31,0.60,0.02}{#1}}
\newcommand{\StringTok}[1]{\textcolor[rgb]{0.31,0.60,0.02}{#1}}
\newcommand{\VariableTok}[1]{\textcolor[rgb]{0.00,0.00,0.00}{#1}}
\newcommand{\VerbatimStringTok}[1]{\textcolor[rgb]{0.31,0.60,0.02}{#1}}
\newcommand{\WarningTok}[1]{\textcolor[rgb]{0.56,0.35,0.01}{\textbf{\textit{#1}}}}
\usepackage{graphicx}
\makeatletter
\newsavebox\pandoc@box
\newcommand*\pandocbounded[1]{% scales image to fit in text height/width
  \sbox\pandoc@box{#1}%
  \Gscale@div\@tempa{\textheight}{\dimexpr\ht\pandoc@box+\dp\pandoc@box\relax}%
  \Gscale@div\@tempb{\linewidth}{\wd\pandoc@box}%
  \ifdim\@tempb\p@<\@tempa\p@\let\@tempa\@tempb\fi% select the smaller of both
  \ifdim\@tempa\p@<\p@\scalebox{\@tempa}{\usebox\pandoc@box}%
  \else\usebox{\pandoc@box}%
  \fi%
}
% Set default figure placement to htbp
\def\fps@figure{htbp}
\makeatother
\setlength{\emergencystretch}{3em} % prevent overfull lines
\providecommand{\tightlist}{%
  \setlength{\itemsep}{0pt}\setlength{\parskip}{0pt}}
\usepackage{bookmark}
\IfFileExists{xurl.sty}{\usepackage{xurl}}{} % add URL line breaks if available
\urlstyle{same}
% fallback for those not using the hyperref driver hyperxmp:
\makeatletter
\@ifundefined{xmpquote}{\newcommand{\xmpquote}[1]{#1}}{}
\makeatother
\hypersetup{
  pdftitle={RPractical\_Montenegro},
  pdfauthor={Jeremy A. Montenegro},
  hidelinks,
  pdfcreator={LaTeX via pandoc}}

\title{RPractical\_Montenegro}
\author{Jeremy A. Montenegro}
\date{2026-11-10}

\begin{document}
\maketitle

Part A: Basic Operation

A1 create v as the sequence 5, 10, 15\ldots50 using seq. print its
length, sum and mean.

\begin{Shaded}
\begin{Highlighting}[]
\NormalTok{v }\OtherTok{\textless{}{-}} \FunctionTok{seq}\NormalTok{(}\DecValTok{5}\NormalTok{, }\AttributeTok{to =} \DecValTok{50}\NormalTok{, }\AttributeTok{by =} \DecValTok{5}\NormalTok{)}
\NormalTok{v}
\end{Highlighting}
\end{Shaded}

\begin{verbatim}
##  [1]  5 10 15 20 25 30 35 40 45 50
\end{verbatim}

\begin{Shaded}
\begin{Highlighting}[]
\FunctionTok{length}\NormalTok{(v)}
\end{Highlighting}
\end{Shaded}

\begin{verbatim}
## [1] 10
\end{verbatim}

\begin{Shaded}
\begin{Highlighting}[]
\FunctionTok{sum}\NormalTok{(v)}
\end{Highlighting}
\end{Shaded}

\begin{verbatim}
## [1] 275
\end{verbatim}

\begin{Shaded}
\begin{Highlighting}[]
\FunctionTok{mean}\NormalTok{(v)}
\end{Highlighting}
\end{Shaded}

\begin{verbatim}
## [1] 27.5
\end{verbatim}

A2 create w that repeats the vector c(1, 2, 3) four times, and z that
repeats each element of c(1, 2, 3) four times. Print both

\begin{Shaded}
\begin{Highlighting}[]
\NormalTok{w }\OtherTok{\textless{}{-}} \FunctionTok{rep}\NormalTok{(}\FunctionTok{c}\NormalTok{(}\DecValTok{1}\NormalTok{, }\DecValTok{2}\NormalTok{, }\DecValTok{3}\NormalTok{), }\AttributeTok{times =} \DecValTok{4}\NormalTok{)}
\NormalTok{z }\OtherTok{\textless{}{-}} \FunctionTok{rep}\NormalTok{(}\FunctionTok{c}\NormalTok{(}\DecValTok{1}\NormalTok{, }\DecValTok{2}\NormalTok{, }\DecValTok{3}\NormalTok{), }\AttributeTok{each =} \DecValTok{4}\NormalTok{)}
\NormalTok{w}
\end{Highlighting}
\end{Shaded}

\begin{verbatim}
##  [1] 1 2 3 1 2 3 1 2 3 1 2 3
\end{verbatim}

\begin{Shaded}
\begin{Highlighting}[]
\NormalTok{z}
\end{Highlighting}
\end{Shaded}

\begin{verbatim}
##  [1] 1 1 1 1 2 2 2 2 3 3 3 3
\end{verbatim}

A3 Compute and print: the remainder of 17 divided by 5, the integer
division of 17 by 5, 2 to the power of 10, the square root of 144, and
pi rounded to 3 decimals.

\begin{Shaded}
\begin{Highlighting}[]
\NormalTok{a }\OtherTok{\textless{}{-}} \DecValTok{17}
\NormalTok{b }\OtherTok{\textless{}{-}} \DecValTok{5}
\NormalTok{divd }\OtherTok{\textless{}{-}}\NormalTok{ a }\SpecialCharTok{/}\NormalTok{ b}
\NormalTok{divd}
\end{Highlighting}
\end{Shaded}

\begin{verbatim}
## [1] 3.4
\end{verbatim}

\begin{Shaded}
\begin{Highlighting}[]
\NormalTok{power }\OtherTok{\textless{}{-}} \DecValTok{2} \SpecialCharTok{\^{}} \DecValTok{10}
\NormalTok{power}
\end{Highlighting}
\end{Shaded}

\begin{verbatim}
## [1] 1024
\end{verbatim}

\begin{Shaded}
\begin{Highlighting}[]
\FunctionTok{sqrt}\NormalTok{(}\DecValTok{144}\NormalTok{)}
\end{Highlighting}
\end{Shaded}

\begin{verbatim}
## [1] 12
\end{verbatim}

\begin{Shaded}
\begin{Highlighting}[]
\NormalTok{w }\OtherTok{\textless{}{-}} \FunctionTok{round}\NormalTok{(pi, }\DecValTok{3}\NormalTok{)}
\NormalTok{w}
\end{Highlighting}
\end{Shaded}

\begin{verbatim}
## [1] 3.142
\end{verbatim}

A4 given a \textless- c(2, 4, 6, 8) and b \textless- c(1, 3, 5, 7),
print a + b, a * b, the logical comparison a \textgreater{} b, and the
dot product sum(a * b).

\begin{Shaded}
\begin{Highlighting}[]
\NormalTok{a }\OtherTok{\textless{}{-}} \FunctionTok{c}\NormalTok{(}\DecValTok{2}\NormalTok{, }\DecValTok{4}\NormalTok{, }\DecValTok{6}\NormalTok{, }\DecValTok{8}\NormalTok{)}
\NormalTok{b }\OtherTok{\textless{}{-}} \FunctionTok{c}\NormalTok{(}\DecValTok{1}\NormalTok{, }\DecValTok{3}\NormalTok{, }\DecValTok{5}\NormalTok{, }\DecValTok{7}\NormalTok{)}
\NormalTok{a }\SpecialCharTok{+}\NormalTok{ b}
\end{Highlighting}
\end{Shaded}

\begin{verbatim}
## [1]  3  7 11 15
\end{verbatim}

\begin{Shaded}
\begin{Highlighting}[]
\NormalTok{a }\SpecialCharTok{*}\NormalTok{ b}
\end{Highlighting}
\end{Shaded}

\begin{verbatim}
## [1]  2 12 30 56
\end{verbatim}

\begin{Shaded}
\begin{Highlighting}[]
\NormalTok{a }\SpecialCharTok{\textgreater{}}\NormalTok{ b}
\end{Highlighting}
\end{Shaded}

\begin{verbatim}
## [1] TRUE TRUE TRUE TRUE
\end{verbatim}

\begin{Shaded}
\begin{Highlighting}[]
\FunctionTok{sum}\NormalTok{(a }\SpecialCharTok{*}\NormalTok{ b)}
\end{Highlighting}
\end{Shaded}

\begin{verbatim}
## [1] 100
\end{verbatim}

A5 Use paste() to produce ``Student-1'' ``Student-2'' ``Student-3''.
Then identify the data type of 42, 42L, ``42'', 41i, and TRUE. What are
its output?

\begin{Shaded}
\begin{Highlighting}[]
\FunctionTok{paste}\NormalTok{(}\StringTok{"Student"}\NormalTok{, }\DecValTok{1}\SpecialCharTok{:}\DecValTok{3}\NormalTok{, }\AttributeTok{sep =} \StringTok{"{-}"}\NormalTok{)}
\end{Highlighting}
\end{Shaded}

\begin{verbatim}
## [1] "Student-1" "Student-2" "Student-3"
\end{verbatim}

\begin{Shaded}
\begin{Highlighting}[]
\FunctionTok{typeof}\NormalTok{(}\DecValTok{42}\NormalTok{)}
\end{Highlighting}
\end{Shaded}

\begin{verbatim}
## [1] "double"
\end{verbatim}

\begin{Shaded}
\begin{Highlighting}[]
\FunctionTok{typeof}\NormalTok{(}\StringTok{"42"}\NormalTok{)}
\end{Highlighting}
\end{Shaded}

\begin{verbatim}
## [1] "character"
\end{verbatim}

\begin{Shaded}
\begin{Highlighting}[]
\FunctionTok{typeof}\NormalTok{(}\DecValTok{41}\DataTypeTok{i}\NormalTok{)}
\end{Highlighting}
\end{Shaded}

\begin{verbatim}
## [1] "complex"
\end{verbatim}

\begin{Shaded}
\begin{Highlighting}[]
\FunctionTok{typeof}\NormalTok{(}\ConstantTok{TRUE}\NormalTok{)}
\end{Highlighting}
\end{Shaded}

\begin{verbatim}
## [1] "logical"
\end{verbatim}

Part B: Data Frames and Conditional Statements

B1 Create a data frame students with 8 rows and these columns (use
stringAsFactors = FALSE; sex is ``M'' or ``F'')

\begin{Shaded}
\begin{Highlighting}[]
\NormalTok{students }\OtherTok{\textless{}{-}} \FunctionTok{data.frame}\NormalTok{(}\AttributeTok{id =} \FunctionTok{seq}\NormalTok{(}\DecValTok{1}\SpecialCharTok{:}\DecValTok{8}\NormalTok{),}
\AttributeTok{name =} \FunctionTok{c}\NormalTok{(}\StringTok{"Alice"}\NormalTok{, }\StringTok{"Bruno"}\NormalTok{, }\StringTok{"Chen"}\NormalTok{, }\StringTok{"Dana"}\NormalTok{, }\StringTok{"Elif"}\NormalTok{, }\StringTok{"Farid"}\NormalTok{, }\StringTok{"Grace"}\NormalTok{, }\StringTok{"Hiro"}\NormalTok{), }
\AttributeTok{sex =} \FunctionTok{c}\NormalTok{(}\StringTok{"F"}\NormalTok{, }\StringTok{"M"}\NormalTok{, }\StringTok{"M"}\NormalTok{, }\StringTok{"F"}\NormalTok{, }\StringTok{"F"}\NormalTok{, }\StringTok{"M"}\NormalTok{, }\StringTok{"F"}\NormalTok{, }\StringTok{"M"}\NormalTok{),}
\AttributeTok{section =} \FunctionTok{c}\NormalTok{(}\StringTok{"A"}\NormalTok{, }\StringTok{"B"}\NormalTok{, }\StringTok{"A"}\NormalTok{, }\StringTok{"C"}\NormalTok{, }\StringTok{"B"}\NormalTok{, }\StringTok{"C"}\NormalTok{, }\StringTok{"A"}\NormalTok{, }\StringTok{"B"}\NormalTok{),}
\AttributeTok{midterm =} \FunctionTok{c}\NormalTok{(}\DecValTok{78}\NormalTok{, }\DecValTok{85}\NormalTok{, }\DecValTok{92}\NormalTok{, }\DecValTok{67}\NormalTok{, }\DecValTok{74}\NormalTok{, }\DecValTok{88}\NormalTok{, }\DecValTok{95}\NormalTok{, }\DecValTok{59}\NormalTok{),}
\AttributeTok{final =} \FunctionTok{c}\NormalTok{(}\DecValTok{82}\NormalTok{, }\DecValTok{79}\NormalTok{, }\DecValTok{96}\NormalTok{, }\DecValTok{71}\NormalTok{, }\DecValTok{80}\NormalTok{, }\DecValTok{91}\NormalTok{, }\DecValTok{93}\NormalTok{, }\DecValTok{65}\NormalTok{),}
\AttributeTok{attendance =} \FunctionTok{c}\NormalTok{(}\DecValTok{95}\NormalTok{, }\DecValTok{88}\NormalTok{, }\DecValTok{100}\NormalTok{, }\DecValTok{72}\NormalTok{, }\DecValTok{90}\NormalTok{, }\DecValTok{85}\NormalTok{, }\DecValTok{98}\NormalTok{, }\DecValTok{60}\NormalTok{))}
\NormalTok{students}
\end{Highlighting}
\end{Shaded}

\begin{verbatim}
##   id  name sex section midterm final attendance
## 1  1 Alice   F       A      78    82         95
## 2  2 Bruno   M       B      85    79         88
## 3  3  Chen   M       A      92    96        100
## 4  4  Dana   F       C      67    71         72
## 5  5  Elif   F       B      74    80         90
## 6  6 Farid   M       C      88    91         85
## 7  7 Grace   F       A      95    93         98
## 8  8  Hiro   M       B      59    65         60
\end{verbatim}

B2 Display the structure, dimension, column names, and summary of
students.

\begin{Shaded}
\begin{Highlighting}[]
\FunctionTok{str}\NormalTok{(students)}
\end{Highlighting}
\end{Shaded}

\begin{verbatim}
## 'data.frame':    8 obs. of  7 variables:
##  $ id        : int  1 2 3 4 5 6 7 8
##  $ name      : chr  "Alice" "Bruno" "Chen" "Dana" ...
##  $ sex       : chr  "F" "M" "M" "F" ...
##  $ section   : chr  "A" "B" "A" "C" ...
##  $ midterm   : num  78 85 92 67 74 88 95 59
##  $ final     : num  82 79 96 71 80 91 93 65
##  $ attendance: num  95 88 100 72 90 85 98 60
\end{verbatim}

\begin{Shaded}
\begin{Highlighting}[]
\FunctionTok{dim}\NormalTok{(students)}
\end{Highlighting}
\end{Shaded}

\begin{verbatim}
## [1] 8 7
\end{verbatim}

\begin{Shaded}
\begin{Highlighting}[]
\FunctionTok{colnames}\NormalTok{(students)}
\end{Highlighting}
\end{Shaded}

\begin{verbatim}
## [1] "id"         "name"       "sex"        "section"    "midterm"   
## [6] "final"      "attendance"
\end{verbatim}

\begin{Shaded}
\begin{Highlighting}[]
\FunctionTok{summary}\NormalTok{(students)}
\end{Highlighting}
\end{Shaded}

\begin{verbatim}
##        id              name          sex         section     midterm     
##  Min.   :1.00   Length   :8   Length   :8   Length   :8   Min.   :59.00  
##  1st Qu.:2.75   N.unique :8   N.unique :2   N.unique :3   1st Qu.:72.25  
##  Median :4.50   N.blank  :0   N.blank  :0   N.blank  :0   Median :81.50  
##  Mean   :4.50   Min.nchar:4   Min.nchar:1   Min.nchar:1   Mean   :79.75  
##  3rd Qu.:6.25   Max.nchar:5   Max.nchar:1   Max.nchar:1   3rd Qu.:89.00  
##  Max.   :8.00                                             Max.   :95.00  
##      final         attendance    
##  Min.   :65.00   Min.   : 60.00  
##  1st Qu.:77.00   1st Qu.: 81.75  
##  Median :81.00   Median : 89.00  
##  Mean   :82.12   Mean   : 86.00  
##  3rd Qu.:91.50   3rd Qu.: 95.75  
##  Max.   :96.00   Max.   :100.00
\end{verbatim}

B3 Add a column average = the mean of midterm and final for each
student, rounded to 1 decimal.

\begin{Shaded}
\begin{Highlighting}[]
\NormalTok{students}\SpecialCharTok{$}\NormalTok{average }\OtherTok{\textless{}{-}} \FunctionTok{round}\NormalTok{((students}\SpecialCharTok{$}\NormalTok{midterm }\SpecialCharTok{+}\NormalTok{ students}\SpecialCharTok{$}\NormalTok{final) }\SpecialCharTok{/} \DecValTok{2}\NormalTok{, }\DecValTok{1}\NormalTok{)}
\NormalTok{students}
\end{Highlighting}
\end{Shaded}

\begin{verbatim}
##   id  name sex section midterm final attendance average
## 1  1 Alice   F       A      78    82         95    80.0
## 2  2 Bruno   M       B      85    79         88    82.0
## 3  3  Chen   M       A      92    96        100    94.0
## 4  4  Dana   F       C      67    71         72    69.0
## 5  5  Elif   F       B      74    80         90    77.0
## 6  6 Farid   M       C      88    91         85    89.5
## 7  7 Grace   F       A      95    93         98    94.0
## 8  8  Hiro   M       B      59    65         60    62.0
\end{verbatim}

B4 Add a logical column passed that is TRUE when average \textgreater=
70 and attendance \textgreater= 75

\begin{Shaded}
\begin{Highlighting}[]
\NormalTok{students}\SpecialCharTok{$}\NormalTok{passed }\OtherTok{\textless{}{-}}\NormalTok{ students}\SpecialCharTok{$}\NormalTok{average }\SpecialCharTok{\textgreater{}=} \DecValTok{70} \SpecialCharTok{\&}\NormalTok{ students}\SpecialCharTok{$}\NormalTok{attendance }\SpecialCharTok{\textgreater{}=} \DecValTok{75}
\NormalTok{students}
\end{Highlighting}
\end{Shaded}

\begin{verbatim}
##   id  name sex section midterm final attendance average passed
## 1  1 Alice   F       A      78    82         95    80.0   TRUE
## 2  2 Bruno   M       B      85    79         88    82.0   TRUE
## 3  3  Chen   M       A      92    96        100    94.0   TRUE
## 4  4  Dana   F       C      67    71         72    69.0  FALSE
## 5  5  Elif   F       B      74    80         90    77.0   TRUE
## 6  6 Farid   M       C      88    91         85    89.5   TRUE
## 7  7 Grace   F       A      95    93         98    94.0   TRUE
## 8  8  Hiro   M       B      59    65         60    62.0  FALSE
\end{verbatim}

B5 Using ifelse(), add a column standing with these values based on
average: ``Honors'' if average \textgreater= 90 ``Good'' if average
\textgreater= 75 (and below 90) ``Needs Support'' otherwise

\begin{Shaded}
\begin{Highlighting}[]
\NormalTok{students}\SpecialCharTok{$}\NormalTok{standing }\OtherTok{\textless{}{-}} \FunctionTok{ifelse}\NormalTok{(students}\SpecialCharTok{$}\NormalTok{average }\SpecialCharTok{\textgreater{}=} \DecValTok{90}\NormalTok{, }\StringTok{"Honors"}\NormalTok{, }\FunctionTok{ifelse}\NormalTok{(students}\SpecialCharTok{$}\NormalTok{average }\SpecialCharTok{\textgreater{}=} \DecValTok{75}\NormalTok{, }\StringTok{"Good"}\NormalTok{, }\StringTok{"Need Support"}\NormalTok{))}
\NormalTok{students}
\end{Highlighting}
\end{Shaded}

\begin{verbatim}
##   id  name sex section midterm final attendance average passed     standing
## 1  1 Alice   F       A      78    82         95    80.0   TRUE         Good
## 2  2 Bruno   M       B      85    79         88    82.0   TRUE         Good
## 3  3  Chen   M       A      92    96        100    94.0   TRUE       Honors
## 4  4  Dana   F       C      67    71         72    69.0  FALSE Need Support
## 5  5  Elif   F       B      74    80         90    77.0   TRUE         Good
## 6  6 Farid   M       C      88    91         85    89.5   TRUE         Good
## 7  7 Grace   F       A      95    93         98    94.0   TRUE       Honors
## 8  8  Hiro   M       B      59    65         60    62.0  FALSE Need Support
\end{verbatim}

B6 Write a function classify(score) that uses an if else if else chain
and returns ``Distinction'' for 85 or above, ``Merit'' for 70 to 84,
``Pass'' for 50 to 69, and ``Fail'' below 50. Test it with sapply(c(90,
73, 55, 40), classify) and print the result.

\begin{Shaded}
\begin{Highlighting}[]
\NormalTok{classify }\OtherTok{\textless{}{-}} \ControlFlowTok{function}\NormalTok{(score) \{}
  \ControlFlowTok{if}\NormalTok{ (score }\SpecialCharTok{\textgreater{}=} \DecValTok{85}\NormalTok{) }
\NormalTok{    \{ }\FunctionTok{return}\NormalTok{(}\StringTok{"Distinction"}\NormalTok{)}
\NormalTok{  \} }\ControlFlowTok{else} 
    \ControlFlowTok{if}\NormalTok{ (score }\SpecialCharTok{\textgreater{}=}  \DecValTok{70} \SpecialCharTok{\&}\NormalTok{ score }\SpecialCharTok{\textless{}=} \DecValTok{84}\NormalTok{) \{ }\FunctionTok{return}\NormalTok{(}\StringTok{"Merit"}\NormalTok{) \} }\ControlFlowTok{else} \ControlFlowTok{if}\NormalTok{ (score }\SpecialCharTok{\textgreater{}=} \DecValTok{50} \SpecialCharTok{\&}\NormalTok{ score }\SpecialCharTok{\textless{}=} \DecValTok{69}\NormalTok{) \{ }\FunctionTok{return}\NormalTok{ (}\StringTok{"Pass"}\NormalTok{) \} }\ControlFlowTok{else}\NormalTok{ \{}\FunctionTok{return}\NormalTok{(}\StringTok{"Fail"}\NormalTok{)\}}
\NormalTok{\}}
\NormalTok{classify}
\end{Highlighting}
\end{Shaded}

\begin{verbatim}
## function (score) 
## {
##     if (score >= 85) {
##         return("Distinction")
##     }
##     else if (score >= 70 & score <= 84) {
##         return("Merit")
##     }
##     else if (score >= 50 & score <= 69) {
##         return("Pass")
##     }
##     else {
##         return("Fail")
##     }
## }
\end{verbatim}

Part C: rbind() and cbind()

C1 Create a data frame new\_students with two students. Its columns must
match the original, including average , passed, and standing, computed
the same way:

id 9, Ines, F, section C, midterm 81, final 84, attendance 92 id 10,
Jamal, M, section A, midterm 70, final 68, attendance 80 Use rbind() to
append them to students, saving the result as students\_all.

\begin{Shaded}
\begin{Highlighting}[]
\NormalTok{new\_students }\OtherTok{\textless{}{-}} \FunctionTok{data.frame}\NormalTok{(}\AttributeTok{id =} \FunctionTok{seq}\NormalTok{(}\DecValTok{9}\SpecialCharTok{:}\DecValTok{10}\NormalTok{),}
\AttributeTok{name =} \FunctionTok{c}\NormalTok{(}\StringTok{"Ines"}\NormalTok{, }\StringTok{"Jamal"}\NormalTok{), }
\AttributeTok{sex =} \FunctionTok{c}\NormalTok{(}\StringTok{"F"}\NormalTok{, }\StringTok{"M"}\NormalTok{), }
\AttributeTok{section =} \FunctionTok{c}\NormalTok{(}\StringTok{"C"}\NormalTok{, }\StringTok{"A"}\NormalTok{), }
\AttributeTok{midterm =} \FunctionTok{c}\NormalTok{(}\DecValTok{81}\NormalTok{, }\DecValTok{70}\NormalTok{), }
\AttributeTok{final =} \FunctionTok{c}\NormalTok{(}\DecValTok{84}\NormalTok{, }\DecValTok{64}\NormalTok{), }
\AttributeTok{attendace =} \FunctionTok{c}\NormalTok{(}\DecValTok{92}\NormalTok{, }\DecValTok{80}\NormalTok{)  }
\NormalTok{)}
\NormalTok{new\_students}
\end{Highlighting}
\end{Shaded}

\begin{verbatim}
##   id  name sex section midterm final attendace
## 1  1  Ines   F       C      81    84        92
## 2  2 Jamal   M       A      70    64        80
\end{verbatim}

C2 Createb the numeric vector quiz \textless- c(8, 7, 10, 6, 7, 9, 10,
5, 8, 6). Use cbind() to attach it to students\_all as a column named
quiz. Confirm the result is still a data frame with class()

C3 Create the matrix m \textless- cbind(midterm =
students\_all\(midterm, final = students_all\)final) and use rbind() to
append a final row containing the column means, labelled ``mean''.

C4 In a comment (2-3 sentence), explain what error rbind() produces if
the data frames have different column names, and show the output if you
safely catch it with tryCatch()

\end{document}
